\section{Conclusion}
\label{sec:conclusion}

This report covered the full arc of a memory-augmented incident-triage project end to end. We built the dataset ourselves --- forking the Online Boutique microservices demo, adding production-realistic OpenTelemetry instrumentation under a strict no-label-leakage rule, defining 27 incident scenario families that cover hard outages, dependency degradation, pod churn, capacity pressure, near-misses, and network failures, running 100 independent fault-injection experiments to produce 6{,}720 labeled telemetry windows, and humanizing a 347-ticket Jira memory corpus in engineer voice (anchored to length and code-block distributions measured directly on the TAWOS public Jira dataset) plus a 110-ticket distractor pool sourced from real TAWOS tickets, fake-incident synthesis in our own services, and cross-architecture infrastructure incidents that we never injected.

We trained seven independent diagnostic pipelines, each representing a different bet about which signal in the telemetry matters most: a gradient-boosted classifier on 94 numeric features (the saturating triage baseline), a fine-tuned dense sentence encoder using contrastive learning on hard-mined negatives (the strongest single retriever for top-1), two variants of a Reciprocal-Rank-Fusion hybrid that combine sparse BM25, dense bi-encoder, and graph-based retrieval (one using regex-extracted entities, one using LLM-extracted entities), a small log-sequence transformer that catches characteristic log patterns the text retrievers miss, a Cypher-traversal retriever over a Neo4j knowledge graph populated by Qwen 3.6 35B-A3B over the 347 humanized tickets, and a three-stage LLM agent (hypothesize, retrieve, verify) that uses thinking-mode reasoning to flag genuinely novel failure modes. Every pipeline won on at least one metric; no pipeline dominated any other.

The cascade --- the Tiered Cascade Hybrid --- composes the seven pipelines into a single four-layer system. The first layer stacks the six triage scores via class-balanced logistic regression and produces a calibrated triage probability. The second layer fuses four retrievers via Reciprocal Rank Fusion, with a separate multi-retriever overlap-rerank rule for the top-1 pick. The third layer detects novel failure modes via three OR-combined signals: the agent's verification output, a retrieval-confidence threshold, and a learned per-window logistic regression over 46 features (window type, scenario family, service name, hard-case flag, max retrieval confidence, triage score). The fourth layer composes the final output. The cascade runs at 16 microseconds per window at inference.

On 1{,}008 held-out test windows the cascade achieves Hit@5 $= 0.912$ (the gold past ticket appears in the top-5 91.2\% of the time), Hit@1 $= 0.722$, MRR $= 0.794$, strict PR-AUC $= 0.9998$, novel-precision $= 0.940$, and novel-recall $= 0.793$. The novel-recall number represents a $+388\%$ relative improvement over the previously locked baseline of $0.163$ at unchanged precision --- the single largest improvement of the project and a result we found by accident more than by design. The original L3 novelty trigger had been wasting easy signal that a 46-feature logistic regression replacing a single fixed threshold could extract for free.

The cascade is robust to memory noise (Hit@5 drops only $2\%$ relative at $50\%$ distractor ratio) and to family-distribution shift (novelty F1 within $11\%$ relative of in-distribution under leave-one-family-out cross-validation). The headline survives both forms of stress.

Three findings transcend any single experiment. \textbf{First}, the cascade is composition-fragile, not component-fragile: a $+404\%$ relative standalone gain on one component broke cascade integration through specific score-scale assumptions. Future cascade extensions must address score normalization explicitly. \textbf{Second}, three independent reranker designs (a fine-tuned cross-encoder, an LLM-agent rerank attempt, and an LLM-judge reranker) all failed in sequence. Multi-retriever consensus on top-1 is empirically dominant on this dataset; no single-LLM judgment from one window can replicate what the cascade's overlap rule extracts from four retrievers' agreement. \textbf{Third}, the retrieval channel was already at $93.5\%$ of its theoretical union ceiling. The real win came from an under-engineered novelty channel, not from retrieval refinement.

The project's contribution to the field is twofold. A reproducible instrumented dataset paired with the construction pipeline that produced it bridges a real gap between public observability datasets (which lack Jira tickets) and public Jira corpora (which lack co-collected telemetry). And a concrete worked example of a four-layer cascade composing seven Pareto-incomparable pipelines, with documented evidence of where the composition holds and where it breaks, gives future practitioners both a design template and a list of pitfalls to avoid.

The path to ICSE submission has one remaining task: finishing the cross-application validation on the OpenTelemetry Demo. The local pilot has cleared all of its gates; the schema-compatibility check passed exactly (94 of 94 feature columns); a clean cross-application finding is already in hand (12 of those 94 features populate from generic OTel telemetry, the remaining 82 depend on Online-Boutique-specific metric names --- the boundary of ``where source-application instrumentation matters''). The full GCP collection is in progress and the two-column zero-shot / L1-retrained external-validity headline is the final shape of the result this paper will report.
